% Editable source for Madgraph_step_User_Guide.pdf.
% Rebuild here with latexmk -pdf -interaction=nonstopmode -halt-on-error Madgraph_step_User_Guide.tex
\documentclass[11pt,letterpaper]{article}
\usepackage[T1]{fontenc}
\usepackage{lmodern}
\usepackage[margin=0.78in,top=0.68in,bottom=0.67in,footskip=24pt]{geometry}
\usepackage{microtype}
\usepackage{xcolor}
\usepackage{array,booktabs,tabularx}
\usepackage{titlesec}
\usepackage{listings}
\usepackage{fancyhdr}
\usepackage[hidelinks]{hyperref}

\definecolor{ink}{HTML}{263846}
\definecolor{navy}{HTML}{193348}
\definecolor{muted}{HTML}{526879}
\definecolor{teal}{HTML}{147078}
\definecolor{pale}{HTML}{F0F5F7}
\definecolor{line}{HTML}{D5E1E7}

\renewcommand{\familydefault}{\sfdefault}
\color{ink}
\setlength{\parindent}{0pt}
\setlength{\parskip}{4pt}
\renewcommand{\arraystretch}{1.21}
\DisableLigatures{encoding = T1, family = lmtt}
\newcommand{\code}[1]{\texttt{\detokenize{#1}}}
\newcolumntype{Y}{>{\raggedright\arraybackslash}X}
\titleformat{\section}[hang]{\large\bfseries\color{navy}}{\color{teal}\thesection.}{0.45em}{}
\titlespacing*{\section}{0pt}{12pt plus 2pt}{4pt}
\hypersetup{pdftitle={Madgraph\_step User Guide},pdfsubject={Terminal workflow and saved results}}
\pagestyle{fancy}
\fancyhf{}
\fancyfoot[L]{\footnotesize\color{muted}Madgraph\_step User Guide}
\fancyfoot[R]{\footnotesize\color{muted}\thepage}
\renewcommand{\headrulewidth}{0pt}
\lstset{
  basicstyle=\ttfamily\small\color{navy},
  backgroundcolor=\color{pale},rulecolor=\color{line},
  frame=single,framerule=0.4pt,framesep=5pt,xleftmargin=5pt,xrightmargin=5pt,
  aboveskip=7pt,belowskip=6pt,columns=fullflexible,keepspaces=true,
  showstringspaces=false,upquote=true
}

\begin{document}

{\Huge\bfseries\color{navy}Madgraph\_step\par}
\vspace{3pt}
{\large\color{muted}A step-by-step terminal guide\par}
\vspace{6pt}
To use your existing \code{input/my_model.txt}, set up the tools in step~1,
then go straight to step~4.

\section{Activate the environment and check the tools}

\begin{lstlisting}
cd "/path/to/Madgraph_step"
conda activate projAI
./run.sh configure-tools
./run.sh check
\end{lstlisting}
Replace the example path with your package folder. Run \code{configure-tools}
on first use or when tool paths change. Review the proposed paths and type
\code{y} to save them to \code{.config}, or \code{edit} to correct them.
The package uses Python 3.11 with NumPy and SciPy, a licensed Mathematica
kernel and the tools listed by \code{check}.

\section{Define particles and generate the model}

Edit \code{input/particles.txt}. Set \code{X} and optionally \code{Y}.
Use \code{Y = None} for a single new multiplet. The file includes
representation examples.

\begin{lstlisting}
./run.sh init input/particles.txt
\end{lstlisting}
This generates \code{input/my_model.txt} with the allowed operators and
starts with diagonal kinetic and mass terms enabled.
\textbf{Running \code{init} again replaces the destination and its edits without a backup.}

\section{Choose operators and coefficients}

In \code{input/my_model.txt}, set \code{enabled} to \code{True} for the operators
you want and supply their coefficients. Edit \code{stable_particles} only for
states you intend to treat as stable. The package derives masses and mixings
using a 125 GeV Higgs and zero VEVs for new scalars.
Use this command to check coefficient units and array order.

\begin{lstlisting}
./run.sh inputs input/my_model.txt
\end{lstlisting}

\section{Validate, list channels and calculate rates}

\begin{lstlisting}
./run.sh validate input/my_model.txt
./run.sh channels input/my_model.txt
./run.sh run input/my_model.txt --selected C0001
\end{lstlisting}
\code{validate} checks the model and derives its spectrum. \code{channels}
exports or reuses the UFO and lists the allowed production channels.
Replace \code{C0001} with an ID from that list. To integrate all listed
production channels, omit \code{--selected}.
The run calculates widths, builds decay histories and saves the ranked rates.

\newpage

{\LARGE\bfseries\color{navy}Find and manage your results\par}

\section{Open the run folder}

Copy the exact path printed next to \textbf{Run files}. With
\lstinline|output = '../output'| in a model inside \code{input/}, the layout is
\code{output/<model_name>/<run-id>/}.

\begin{tabularx}{\linewidth}{@{}p{0.39\linewidth}Y@{}}
\toprule
\textbf{File} & \textbf{What to look for} \\
\midrule
\code{report.html} & Open in a browser for the run summary and links. \\
\code{ranking.csv} & Decay histories ranked by production cross section $\times$ cascade branching probability, in pb. \\
\code{production_channels.csv} & The available production channels and their IDs. \\
\code{cross_sections.csv} & Integrated production channels and their numerical errors. \\
\code{total_widths.csv}, \code{decays.csv} & Total widths, decay modes and branching fractions. \\
\code{manifest.json} & Completion status. Check \code{requested_complete} for the selected run. \\
\bottomrule
\end{tabularx}

Open \code{report.html} to see the mass spectrum and enabled couplings, then production cross sections, widths and ranked decay histories. Final-state filters apply to the displayed page. Choose \textbf{Download ranked rates (CSV)} for all saved rows. To refresh a saved run, including one moved to another directory, use its run folder:
\begin{lstlisting}
./run.sh report --run-dir "path/to/saved run"
\end{lstlisting}
This reads saved results and writes HTML without running Mathematica or MadGraph.

A blank rate means its production channel has no completed integration.
These are tree-level, narrow-width estimates. Keep overlapping histories
separate and check the report for omitted decay contributions.

\section*{Use a different input filename}

For a \textbf{new model}, choose the generated filename with \code{--model}.
You can also pass a renamed particle source in place of \code{input/particles.txt}.

\begin{lstlisting}
./run.sh init input/particles.txt --model "input/my scan.txt"
\end{lstlisting}
To \textbf{keep an existing model's edits}, rename the TXT file directly.
Leave its \code{_metadata} line and the adjacent \code{.model_metadata/}
folder in place. If moving to another folder, bring the metadata folder too.

\begin{lstlisting}
mv -i input/my_model.txt "input/my scan.txt"
./run.sh channels "input/my scan.txt"
\end{lstlisting}
Pass the chosen filename to each model command, or edit \code{model = input/my scan.txt}
under \code{[project]} in \code{.config} to make it the default.
To name the results subfolder, set \code{model_name} inside the TXT file.
A relative \code{output} in that file starts beside the TXT file.
Use \code{--output "scan results"} on a run command to override it for that run.
A relative path in this option starts in the terminal's current folder.

\section*{Inspect, stop or resume}

\begin{lstlisting}
./run.sh decays input/my_model.txt
./run.sh histories input/my_model.txt --page-size 20
\end{lstlisting}
\code{decays} calculates widths and branching fractions. \code{histories}
also builds the cascade catalogue without integrating production.
\code{--page-size} limits the display only. A selected production run may
still build histories for all discovered channels. To restrict production
to at most two outgoing new particles, set \code{max_final} to \code{2}
and \code{new_particles_only} to \code{True} in \code{collider}.
These settings apply before decay.

Press \textbf{Ctrl+C} to stop and repeat the same command to resume from
verified checkpoints. If the package flags stale coefficient bindings,
run \code{./run.sh setup input/my_model.txt} and review the values it requests.

\end{document}
